% Paul E. West and Sean T. Hayes

\documentclass[14pt,aspectratio=1610, hyperref={pdfpagelayout=SinglePage}]{beamer}
% \documentclass[14pt,aspectratio=1610, hyperref={pdfpagelayout=SinglePage},handout]{beamer}
\usepackage{csu-theme}
\usepackage{booktabs}

\title{\Huge{}Operator\\ Overloading}
\subtitle{CSCI 315: \textit{Data Structure Analysis}\\Chapter 13}
%\author{Sean T. Hayes, Ph.D.}
\institute{Charleston Southern University}

%% Comment out date to use default (today)
% \date{January 30, 2025}
\subject{Software Programming}
%\keywords{}

\setlength{\parskip}{0.5em}

\begin{document}

\begin{frame}
  \titlepage%
\end{frame}

\section{Motivation}

\begin{frame}[fragile]{Objectives on Opeartor Overloading}

In this lecture, you will:
\begin{itemize}
\item
  Examine the purpose and limitations of operator overloading.
\item
  Examine the \lstinline{this} pointer.
\item
  Differentiate between member and nonmember (\lstinline{friend}) operator overloads.
\item
  Overload operators as members and nonmembers of a class.
\item
  Apply the \emph{Rule of Three} to properly handle pointer member variables.
\end{itemize}
\end{frame}

\begin{frame}[fragile]{Why Operator Overloading Is Needed}
\protect\hypertarget{why-operator-overloading-is-needed-1-of-2}{}

Consider the following statements:
\begin{lstlisting}
ClockType local{8, 23, 34};
ClockType yourTime{4, 23, 30};
\end{lstlisting}

Which version of the C++ statements below would you prefer?
\begin{columns}[T]
\column{0.5\textwidth}
\begin{lstlisting}[basicstyle=\ttfamily\small]
local.print();
local.incrementSeconds();
if (local.equal(yourTime))
// ...
\end{lstlisting}

\column{0.43\textwidth}
\begin{lstlisting}[basicstyle=\ttfamily\small, xleftmargin=0pt]
cout << local;
local++;
if (local == yourTime)
// ...
\end{lstlisting}
\end{columns}
\end{frame}

\begin{frame}{Why Operator Overloading Is Needed}
\protect\hypertarget{why-operator-overloading-is-needed-2-of-2}{}
\begin{itemize}
\item
  \textit{Assignment} and \textit{member selection} are the only \\built-in operations on
  classes.
  \begin{itemize}
  \item
    Other operators cannot be applied directly to class objects.
  \end{itemize}
\item
  \emph{Operator overloading} extends the definition of an operator to
  work with a user-defined data type.
  \begin{itemize}
  \item
    C++ allows you to extend the definitions of most of the operators to
    work with classes.
  \end{itemize}
\end{itemize}
\end{frame}

\section{Operator Overloading}

\begin{frame}[fragile]{Operator Overloading}
\protect\hypertarget{operator-overloading}{}
\begin{itemize}
\item
  Most existing C++ operators can be overloaded to manipulate class
  objects.
\item
  New operators cannot be created.
\item
  An \emph{operator function} is a function that overloads an operator.
  \begin{itemize}
  \item
    Use reserved word \lstinline{operator} followed by the operator as
    the function name.
  \end{itemize}
\end{itemize}
\end{frame}

\begin{frame}[fragile]{Syntax for Operator Functions}
\protect\hypertarget{syntax-for-operator-functions}{}
\begin{itemize}
\item
  Syntax of an operator function heading:\\
  \lstinline{returnType operator symbol(formal parameter list)}
  \begin{itemize}
  \item
    It is a value-returning function.
  \item
    \lstinline{operator} is a reserved word.
  \end{itemize}
\item
  To overload an operator for a class:
  \begin{itemize}
  \item
    Include the operator function declaration in the class definition.
  \item
    Write the definition of the operator function.
  \end{itemize}
\end{itemize}
\end{frame}



\begin{frame}[fragile]{Friend Functions of Classes}
\protect\hypertarget{friend-functions-of-classes}{}
\begin{itemize}
\item
  A \emph{friend function} (of a class) is a nonmember function of the class 
  that has access to all the members of the class.
\item
  Use the reserved word \lstinline{friend} in the function prototype in the
  class definition.
\item
  Friendship is always given by the class.

\begin{lstlisting}[basicstyle=\ttfamily\small]
class IllusFriend
{
    friend void two(/*parameters*/);
    // ...
};
\end{lstlisting}
\end{itemize}
\end{frame}

\begin{frame}[fragile]{Definition of a Friend Function}
\protect\hypertarget{definition-of-a-friend-function}{}
\begin{itemize}
\item
  \lstinline{friend} does not appear in the heading of the function's
  definition.
\item
  When writing the \lstinline{friend} function's definition.
\item
  The name of the class and the scope resolution operator are not used.

\begin{lstlisting}[basicstyle=\ttfamily\small]
void two(/*parameters*/)
{
    // ...
}
\end{lstlisting}
\end{itemize}
\end{frame}

\begin{frame}[fragile]{Operator Functions as Member and Nonmember
Functions}
\protect\hypertarget{operator-functions-as-member-and-nonmember-functions}{}
\begin{itemize}
\item
  To overload \lstinline{()}, \lstinline{[]},
  \lstinline{->}, or \lstinline{=} for a class,
  the function \textit{must} be a member of the class.
\item
  Suppose \lstinline{op} is overloaded for \lstinline{opOverClass}:
  \begin{itemize}
  \item
    If the leftmost operand of \lstinline{op} is an object of a different
    type, the overloading function must be a nonmember (friend) of the
    class.
  \item
    If the overloading function for \lstinline{op} is a member of
    \lstinline{opOverClass}, then when applying \lstinline{op} on objects of
    type \lstinline{opOverClass}, the leftmost operand must be of type
    \lstinline{opOverClass}.
  \end{itemize}
\end{itemize}
\end{frame}

\section{Binary Operators}

\begin{frame}[fragile]{Overloading Binary Operators}
\protect\hypertarget{overloading-binary-operators}{}
\begin{itemize}
\item
  If \texttt{\textbf{\#}} represents a binary operator (e.g.,
  \lstinline{+} or \lstinline{==}) that is to be overloaded for \lstinline{rectangleType}.
\item
  It can be overloaded as either a member function of the class or as a
  friend function.
\end{itemize}
\end{frame}

\begin{frame}[fragile]{Binary Operator Overloads as Methods}
\protect\hypertarget{overloading-the-binary-operators-as-member-functions}{}
\begin{itemize}
\item
  Function prototype (included in the class definition):\\
  \lstinline{returnType operator#(const Type&) const;}
\item
  Function definition:\\
  \begin{lstlisting}[basicstyle=\ttfamily\small]
returnType ClassType::operator#(const Type& otherObject) const
{
  // algorithm to perform the operation
  return value;
}
\end{lstlisting}
\end{itemize}
\end{frame}

\begin{frame}[fragile]{Overloading the Assignment Operator (\texttt{=})}
\protect\hypertarget{overloading-the-assignment-operator}{}
  Function prototype:
  \lstinline{const Type& operator=(const Type&);}

  Function definition:
  \begin{lstlisting}[basicstyle=\ttfamily\small]
const Type& operator=(const Type& rightObject)
{
  // local declaration, if any
  // Avoid self-assignment
  if (this != &rightObject)
  {
    // Copy rightObject into this object
  }
  return *this; // return the object assigned
}\end{lstlisting}

\end{frame}

\begin{frame}[fragile]{Binary Operator Overloads as Nonmember}
\protect\hypertarget{overloading-the-binary-operators-arithmetic-or-relational-as-nonmember-functions}{}
\begin{itemize}
\item
  Function prototype (included in class definition):
  \begin{lstlisting}
friend returnType operator#(const Type&, const Type&) const;
\end{lstlisting}
\item
  Function definition:
  \begin{lstlisting}[basicstyle=\ttfamily\small]
returnType operator#(const Type& leftObject, const Type& rightObject) const
{
  // algorithm to perform the operation
  return value;
}
\end{lstlisting}
\end{itemize}
\end{frame}

\begin{frame}[fragile]{Insertion
(\texttt{{<}<}) and Extraction (\texttt{{>}>}) Operators}
\begin{itemize}

\item
  Consider the expression:\\
  \lstinline{cout << myRectangle;}

  \begin{itemize}
  \item
    Leftmost operand is an \lstinline{std::ostream} object, not a
    \lstinline{rectangleType} object.
  \end{itemize}
\item
  The operator overload of
  \lstinline{<<} for \lstinline{rectangleType} must be a
  nonmember function of the class.
  \begin{itemize}
  \item
    The same applies to the function that overloads
    \lstinline{>>}.
    \end{itemize}
\end{itemize}
\end{frame}

\begin{frame}[fragile]{Overloading the Stream Insertion Operator
(\texttt{{<}<})}
\protect\hypertarget{overloading-the-stream-insertion-operator}{}

  Function prototype:
  \begin{lstlisting}
friend std::ostream& operator<<( std::ostream&, const Type&);
\end{lstlisting}

  Function definition:
  \begin{lstlisting}
std::ostream& operator<<(std::ostream& out,
      const Type& cObject) {
  // local declaration, if any
  // Output the members of cObject.
  // out << ...

  return out; // return the stream object.
}\end{lstlisting}

\end{frame}

\begin{frame}[fragile]{Overloading the Stream Extraction Operator
  (\texttt{{>}>})}

  Function prototype:
  \begin{lstlisting}
friend std::istream& operator>>( std::istream&, Type&);
\end{lstlisting}

  Function definition:
  \begin{lstlisting}
std::istream& operator>>(std::istream& in, Type& cObject) {
  // local declaration, if any
  // Read the members of cObject.
  // in >> ...

  return in; // return the stream object.
}\end{lstlisting}

\end{frame}

\begin{frame}{Classes with Pointer Members (Revisited)}
\begin{itemize}
\item
  Recall that the assignment operator copies member variables from one
  object to another of the same type.
  \begin{itemize}
  \item
    Does not work well with pointer member variables.
  \end{itemize}
\item
  Classes with pointer member variables must:
  \begin{enumerate}
  \item
    Explicitly overload the assignment operator.
  \item
    Include the copy constructor.
  \item
    Include the destructor.
  \end{enumerate}
\end{itemize}
\end{frame}

\begin{frame}[fragile]{Member vs. Nonmember Operator Overloading}
Some operators must be overloaded as member functions and some as nonmember (\lstinline{friend}) functions.

Binary arithmetic operators can be overloaded either way.
  \begin{lstlisting}
Type operator+(const Type&);
\end{lstlisting}
  or
\begin{lstlisting}
friend Type operator+(const Type&, const Type&);
\end{lstlisting}
\end{frame}

\begin{frame}[fragile]{Member vs. Nonmember Operator Overloading}
Member Function for Operator Overload:
\begin{itemize}
  \item
    As a member function, the operation has direct\\ access to data
    members of one of the objects.
  \item
    Need to pass only one object as a parameter.
\end{itemize}

Non-Member Function for Operator Overload:
\begin{itemize}
  \item
    Both operands are included as parameters when nonmember functions that define binary operators.
  \item
    The code may be somewhat clearer this way.
\end{itemize}
\end{frame}

\begin{frame}[fragile]{Overloading the Array Index Operator \texttt{[]}}
Declaring the \lstinline{[]} operator as a class member for
  nonconstant arrays:
\begin{lstlisting}
Type& operator[](unsigned int index);
\end{lstlisting}

Declaring the \lstinline{[]} operator as a class member for
  constant arrays:
\begin{lstlisting}
const Type& operator[](unsigned int index) const;
\end{lstlisting}
\end{frame}

\begin{frame}[fragile]{Overloading the Array Index Operator \texttt{[]}}
\begin{lstlisting}
int array[5] { 1, 2, 3, 4, 5 };
const int ARRAY[5] { 1, 2, 3, 4, 5 };

std::cout << array[2];
std::cout << ARRAY[2]; // Const operator[]

array[2] = 10;    // Non-const operator[]
// ARRAY[2] = 10; // Error
\end{lstlisting}
\end{frame}

\section{Unary Operator}

\begin{frame}[fragile]{Overloading Unary Operators}
\begin{itemize}
\item
  To overload a unary operator for a class:
  \begin{itemize}
  \item
    If the operator function is a member of the class, it has no
    parameters.
  \item
    If the operator function is a nonmember (i.e., a \lstinline{friend}
    function), it has one parameter.
  \end{itemize}
\end{itemize}
\end{frame}

\begin{frame}[fragile]{Pre-Increment (\texttt{++}) and
Pre-Decrement (\texttt{{-}-})}

The general syntax to overload the pre-increment operator \lstinline{++} as a member function.

\begin{lstlisting}
Type operator++() {
  // Increment the value of the object by 1
  return *this; //return the new value
}
\end{lstlisting}
\end{frame}

\begin{frame}[fragile]{Post-Increment (\texttt{++}) and
Post-Decrement (\texttt{{-}-})}

The general syntax to overload the post-increment operator \lstinline{++} as
a member function.

\begin{lstlisting}
Type operator++(int) {
  const auto TEMP = *this; // copy old value

  // Increment the value of the object by 1

  return TEMP; // the copy of the old value
}
\end{lstlisting}
\end{frame}

\begin{frame}[fragile]{Post-Increment (\texttt{++}) and
Post-Decrement (\texttt{{-}-})}
The general syntax to overload the post-increment operator \lstinline{++} as
  a nonmember function.

Function definition:
\begin{lstlisting}
Type operator++(Type& incObj, int) {
  const auto TEMP = incObj; // copy old

  // Increment the value of the object by 1

  return TEMP; // the copy of the old value
}
\end{lstlisting}
May need to declare as a \lstinline{friend} function in the class definition.
\end{frame}

\section{Restrictions}
\begin{frame}[fragile]{Overloading an Operator: Some Restrictions}

\emph{Only the behavior is customizable},\\ not the syntax rules.

The following are all fixed by the language and\\ cannot be changed:

\pause{}
\begin{itemize}[<+->]
\item
  Precedence

\item
  Associativity

\item
  Arity (unary vs binary)\\
  Cannot change number of parameters.

\item
  Short-circuit evaluation for \lstinline{&&} and \lstinline{||}
\end{itemize}
\end{frame}

\begin{frame}[fragile]{Overloading an Operator: Some Restrictions}
\begin{itemize}[<+->]
\item
  Cannot create new operators.
\item
  Cannot overload: \lstinline{.}~~\lstinline{.*}~~\lstinline{::}~~\lstinline{?:}~~\lstinline{sizeof}~~\\
  and casting operators.
\item
  Operators on built-in types cannot be redefined.\\
  For example, cannot redefine \lstinline{+} for two \lstinline{int} values.
\item
  Can overload for user-defined objects or for a combination of
  user-defined and built-in objects.
\end{itemize}
\end{frame}


% \begin{frame}[fragile]{Operator Overloading: One Final Word}
% \protect\hypertarget{operator-overloading-one-final-word}{}
% \begin{itemize}
% \item
%   If an operator \texttt{op} is overloaded for a class, e.g.,
%   \texttt{rectangleType}
% \item
%   When you use \texttt{op} on objects of type \texttt{rectangleType},
%   the body of the function that overloads the operator \texttt{op} for
%   the \texttt{class\ rectangleType} executes
% \item
%   Therefore, whatever code you put in the body of the function executes
% \end{itemize}
% \end{frame}

% \begin{frame}{Function Overloading}
% \protect\hypertarget{function-overloading}{}
% \begin{itemize}
% \item
%   Overloading a function refers to having several functions with the
%   same name, but different parameters
% \item
%   The parameter list determines which function will execute
% \item
%   Must define each function
% \end{itemize}
% \end{frame}

\section{Review}
\begin{frame}[fragile]{Quick Review}
\begin{itemize}[<+->]
\item
  \emph{Overloaded operator} -- an operator that has different meanings with different data types.
\item
  \emph{Operator function} -- a function that overloads an operator.
\item
  Operator functions are value-returning.
\item
  Operator overloading provides the same concise notation for
  user-defined data types as for built-in data types.
\end{itemize}
\end{frame}

\begin{frame}[fragile]{Quick Review}
\begin{itemize}[<+->]
\item
  Only existing operators can be overloaded.
\item
  The \lstinline{this} pointer refers to the object.
\item
  A \lstinline{friend} function is a nonmember of a class.
\item
  If an operator function is a member of a class,
  the leftmost operand of the operator must be a class object (or a
  reference to a class object) of that operator's class.
\end{itemize}
\end{frame}

\begin{frame}[fragile]{Quick Review}
\begin{itemize}
\item
  \emph{Rule of 3}: Classes with pointer variables\\ \emph{must} overload the
  \pause{}
  \begin{enumerate}[<+->]
  \item
    copy constructor,
  \item 
    destructor, and
  \item 
    assignment operator.
  \end{enumerate}
\end{itemize}
\end{frame}

\end{document}
